\section{局限、失败模式与相关机制}

\subsection{为什么这些方法会卡住}

这一讲的 toy 系统其实已经暴露了策略梯度的几个典型失败模式：
\begin{itemize}
    \item \textbf{reward 稀疏}：大多数 sample 没信号，更新太慢；
    \item \textbf{baseline 不准}：group 太小或方差太大时，centered signal 仍然不稳；
    \item \textbf{更新过猛}：ratio 太大，模型可能一下子偏离原策略；
    \item \textbf{KL 太弱}：模型会为了 reward 牺牲原本的语言能力；
    \item \textbf{KL 太强}：模型几乎不动，reward 学不上去。
\end{itemize}

\begin{table}[H]
\centering
\begin{tabular}{@{}p{3.1cm}p{4.5cm}p{5.1cm}@{}}
\toprule
\textbf{问题} & \textbf{表面现象} & \textbf{本质原因} \\
\midrule
稀疏 reward & 训练曲线长时间不动 & 梯度大多数时候接近 0 \\
group 太小 & centered 后仍很抖 & baseline 估计噪声大 \\
ratio 太大 & loss 突然震荡 & 单步更新过猛 \\
KL 太弱 & 语言质量下降 & 目标函数只看 reward，不看漂移 \\
KL 太强 & reward 学不起来 & 约束压过了优化 \\
\bottomrule
\end{tabular}
\caption{GRPO / PPO 风格训练最常见的失败模式：问题通常不在“公式错了”，而在“信号和约束没平衡好”}
\end{table}

\subsection{一个 worked example: group mean 怎么决定更新方向}

设同一个 prompt 采到三个 response，它们的 reward 是 $[3,2,1]$。那么 group mean 是 2，centered rewards 就是 $[+1,0,-1]$。这意味着：
\begin{itemize}
    \item 第一条 response 被强化；
    \item 第二条 response 基本不动；
    \item 第三条 response 被压低。
\end{itemize}

如果把它换成 normalized rewards，结果只是再除以一个标准差；更新方向不变，主要变化是步幅更统一。


\subsection{GRPO 和相关 alignment 方法的关系}

如果把这门课放到更大的对齐框架里看，GRPO 和 RLHF、PPO、REINFORCE 的关系可以总结成一句话：\textbf{它们共享同一个 policy gradient 骨架，只是在 baseline、约束和奖励来源上做了不同选择}。

\begin{table}[H]
\centering
\begin{tabular}{@{}p{2.7cm}p{3.7cm}p{3.7cm}p{3.5cm}@{}}
\toprule
\textbf{方法} & \textbf{reward 来源} & \textbf{约束} & \textbf{适用感受} \\
\midrule
REINFORCE & 直接 reward & 通常没有显式约束 & 最基础，但方差大 \\
PPO & reward model 或环境 reward & clipping + value baseline & 经典、稳定、但组件多 \\
GRPO & group-relative reward & clipping + KL, 不一定要 critic & 对 LM response group 很自然 \\
RLHF 训练链路 & 人类偏好或奖励模型 & 取决于实现 & 更偏系统工程 \\
\bottomrule
\end{tabular}
\caption{把 GRPO 放到更大的 alignment 图景里：它不是另起炉灶，而是把现有 policy gradient 组件重新组合}
\end{table}

\begin{importantbox}{这一节想传达的判断}
GRPO 的价值不是“发明了一个完全不同的 RL”，而是它把 \emph{group baseline}、\emph{clipping}、\emph{KL} 这些已经成熟的稳定化技巧放进了语言模型最自然的采样结构里，因此实现上更轻，直觉上也更顺。
\end{importantbox}

